\chapter{Fundamental Equations}

This chapter mixes equations from different levels of physics. Start with the conservation laws and Newton's laws, then return to Maxwell's equations, relativity, and quantum mechanics after the symbols feel familiar. For each item, read the sentence before the formula first; the formula is the compact version of that sentence.

\textbf{Reality Mapping:} These equations are the grammar of physics---they tell us what can and cannot happen in the universe. Conservation laws say that certain quantities (energy, momentum, angular momentum) are the same before and after any process. Maxwell's equations say that electric and magnetic fields are two aspects of one thing, and that light is an electromagnetic wave. Einstein's equations say that mass and energy are interchangeable, and that gravity is the curvature of spacetime. Quantum equations say that particles behave like waves, that you cannot know everything at once, and that the vacuum is seething with activity. Together, these equations describe the rules that govern everything from falling apples to exploding stars.

\begin{enumerate}

\item \textbf{Ampere--Maxwell Law:} Electric currents and changing electric fields both create magnetic fields. This equation also predicts that light is an electromagnetic wave.
\[
\nabla\times\mathbf B=\mu_0\mathbf J+\mu_0\epsilon_0\frac{\partial\mathbf E}{\partial t}
\]
\begin{itemize}
\item \textbf{Assumptions:} The medium is linear and time-invariant. The fields vary slowly enough that radiation reaction is negligible.
\item \textbf{Limitations:} Fails at quantum scales (need quantum electrodynamics). Does not account for magnetic monopoles (none have been observed). In media with large dispersion or nonlinear effects, additional terms are needed.
\item \textbf{Constraints:} Valid in the classical regime where $v\ll c$ for material motion. Requires $\mu_0$ and $\epsilon_0$ to be constants (true in vacuum).
\end{itemize}

\item \textbf{Conservation of Angular Momentum:} The total angular momentum of an isolated system does not change. This is why ice skaters spin faster when they pull their arms in.
\[
\mathbf L=\text{const}
\]
\begin{itemize}
\item \textbf{Assumptions:} The system is isolated---no external torques act on it. The Lagrangian of the system is invariant under rotations.
\item \textbf{Limitations:} Does not apply in non-inertial (rotating) frames without adding fictitious torques. In general relativity, angular momentum is conserved only locally.
\item \textbf{Constraints:} Valid for closed systems only. In open systems, external torques must be included.
\end{itemize}

\item \textbf{Conservation of Energy:} Energy can change form but the total amount stays the same.
\[
\Delta E_{\text{total}}=0
\]
\begin{itemize}
\item \textbf{Assumptions:} The system is closed (no energy enters or leaves). Time-translation symmetry holds (the laws of physics do not change with time).
\item \textbf{Limitations:} In an expanding universe, energy conservation becomes subtle (photons lose energy to cosmological redshift). In general relativity, global energy conservation is not well-defined in all cases.
\item \textbf{Constraints:} Applies strictly to isolated systems. In open systems, energy input and output must be accounted for.
\end{itemize}

\item \textbf{Conservation of Mass:} Mass is neither created nor destroyed.
\[
\Delta m=0
\]
\begin{itemize}
\item \textbf{Assumptions:} Velocities are much less than the speed of light ($v\ll c$). No nuclear reactions or particle--antiparticle annihilation occur.
\item \textbf{Limitations:} Fails at relativistic speeds---mass can be converted to energy and vice versa via $E=mc^2$. In nuclear and particle physics, mass is not separately conserved.
\item \textbf{Constraints:} Valid only in classical, non-relativistic mechanics. In chemistry and everyday life, it holds to extraordinary precision.
\end{itemize}

\item \textbf{Conservation of Momentum:} The total momentum of an isolated system does not change.
\[
\sum\mathbf p_i=\text{const}
\]
\begin{itemize}
\item \textbf{Assumptions:} The system is isolated---no external forces act on it. The Lagrangian is invariant under spatial translations.
\item \textbf{Limitations:} In an expanding universe, momentum of photons is not conserved globally. In general relativity, momentum conservation is local, not global.
\item \textbf{Constraints:} Valid for closed systems. In open systems, external impulses must be included.
\end{itemize}

\item \textbf{Einstein Field Equations:} The core of general relativity. They say that mass and energy curve spacetime, and curved spacetime tells mass how to move. They predict black holes, gravitational waves, and the expansion of the universe.
\[
R_{\mu\nu}-\frac12Rg_{\mu\nu}+\Lambda g_{\mu\nu}
=\frac{8\pi G}{c^4}T_{\mu\nu}
\]
\begin{itemize}
\item \textbf{Assumptions:} Spacetime is a smooth, differentiable manifold. The metric tensor fully describes gravity. Matter and energy are described by a stress--energy tensor $T_{\mu\nu}$.
\item \textbf{Limitations:} Does not incorporate quantum effects---breaks down at the Planck scale ($\sim10^{-35}$\,m). Cannot describe the singularity inside a black hole. Quantum gravity is needed for a complete theory.
\item \textbf{Constraints:} Valid in the classical regime. Requires knowledge of the stress--energy tensor, which must be supplied by additional physics (e.g., fluid dynamics, electrodynamics).
\end{itemize}

\item \textbf{Faraday's Law:} A changing magnetic field creates an electric field. This is how electric generators work.
\[
\nabla\times\mathbf E=-\frac{\partial\mathbf B}{\partial t}
\]
\begin{itemize}
\item \textbf{Assumptions:} The electric and magnetic fields vary slowly compared to the speed of light. The medium is linear and time-invariant.
\item \textbf{Limitations:} Does not account for quantum effects (e.g., the Aharonov--Bohm effect). Fails in nonlinear media or at very high field intensities.
\item \textbf{Constraints:} Valid in the classical regime. The minus sign (Lenz's law) assumes no magnetic monopoles.
\end{itemize}

\item \textbf{Gauss's Law:} Electric charges produce electric fields. The total electric flux through a closed surface equals the enclosed charge divided by $\epsilon_0$.
\[
\nabla\cdot\mathbf E=\frac{\rho}{\epsilon_0}
\]
\begin{itemize}
\item \textbf{Assumptions:} The electric field is created by static or slowly varying charge distributions. The permittivity $\epsilon_0$ is a constant (true in vacuum).
\item \textbf{Limitations:} Does not apply in quantum electrodynamics where virtual charges modify the effective vacuum permittivity. In dielectric media, the effective $\epsilon$ replaces $\epsilon_0$.
\item \textbf{Constraints:} Valid for classical electrodynamics. Requires that magnetic monopoles do not exist (consistent with all observations to date).
\end{itemize}

\item \textbf{Gauss's Law for Magnetism:} There are no magnetic monopoles---magnetic field lines always form closed loops.
\[
\nabla\cdot\mathbf B=0
\]
\begin{itemize}
\item \textbf{Assumptions:} Magnetic monopoles do not exist. The divergence of the magnetic field is zero everywhere.
\item \textbf{Limitations:} If magnetic monopoles were ever discovered, this equation would need to be modified (adding a magnetic charge density term). Some grand unified theories predict monopoles, but none have been observed.
\item \textbf{Constraints:} Valid in all observed physical situations. Consistent with Maxwell's equations and all experimental data.
\end{itemize}

\item \textbf{Law of Thermodynamics (First):} Energy is conserved. The heat added to a system goes into increasing its internal energy or doing work.
\[
dU=\delta Q-\delta W
\]
\begin{itemize}
\item \textbf{Assumptions:} The system is closed (no mass enters or leaves). Energy can be exchanged as heat or work, but not created or destroyed.
\item \textbf{Limitations:} Does not predict the direction of processes (that is the second law). In relativistic systems, mass--energy equivalence must be included.
\item \textbf{Constraints:} Applies to any thermodynamic process, reversible or irreversible. The internal energy $U$ is a state function; $\delta Q$ and $\delta W$ are path-dependent.
\end{itemize}

\item \textbf{Law of Thermodynamics (Second):} Entropy (disorder) of an isolated system never decreases. Heat flows from hot to cold, never the other way.
\[
\Delta S\geq 0
\]
\begin{itemize}
\item \textbf{Assumptions:} The system is isolated (no energy or matter exchange with surroundings). The system is in or approaches thermal equilibrium.
\item \textbf{Limitations:} Statistical in nature---entropy can locally decrease with extremely small probability. In small systems (nanoscale), thermal fluctuations can temporarily decrease entropy. Does not apply to non-equilibrium systems without modification.
\item \textbf{Constraints:} Strictly valid only for isolated systems. For open systems, the entropy of the system plus its surroundings must increase.
\end{itemize}

\item \textbf{Law of Thermodynamics (Third):} As temperature approaches absolute zero, entropy approaches a minimum value. You can never actually reach absolute zero.
\[
\lim_{T\to 0}S=S_0
\]
\begin{itemize}
\item \textbf{Assumptions:} The system is in its quantum mechanical ground state. The third law applies to perfect crystals where $S_0=0$. Systems with residual entropy (e.g., ice) have $S_0>0$.
\item \textbf{Limitations:} Does not apply to glasses or disordered systems, which can have residual entropy at $T=0$. The unattainability statement is not rigorously proven for all systems.
\item \textbf{Constraints:} Requires a perfect crystal at absolute zero to have zero entropy. Practically, absolute zero cannot be reached in finite steps (Nernst's unattainability statement).
\end{itemize}

\item \textbf{Law of Thermodynamics (Zeroth):} If two objects are both in thermal equilibrium with a third, they are in equilibrium with each other. This defines temperature.
\[
\text{If }A\sim C\text{ and }B\sim C,\text{ then }A\sim B
\]
\begin{itemize}
\item \textbf{Assumptions:} Thermal equilibrium is a transitive relation. The systems can exchange energy (or have been in contact) and reach a common temperature.
\item \textbf{Limitations:} Does not define temperature in non-equilibrium systems. Breaks down for systems with long-range interactions (e.g., gravitational systems) where thermal equilibrium may not be well-defined.
\item \textbf{Constraints:} Valid for systems in local thermal equilibrium. Assumes that temperature is a well-defined scalar quantity.
\end{itemize}

\item \textbf{Lorentz Transformation:} The equations that replace Newton's Galilean transformations when objects move close to the speed of light. They show that time slows down and lengths contract for moving objects, and they ensure that the speed of light is the same for everyone.
\[
x'=\gamma(x-vt),\qquad
t'=\gamma\left(t-\frac{vx}{c^2}\right)
\]
\begin{itemize}
\item \textbf{Assumptions:} The speed of light $c$ is the same in all inertial reference frames. Space is flat (no gravity). The two frames are related by a constant relative velocity $v$.
\item \textbf{Limitations:} Does not apply in curved spacetime (general relativity uses general coordinate transformations). Breaks down for accelerated observers (though it can be applied locally).
\item \textbf{Constraints:} Valid only for $v<c$. For $v\geq c$, the Lorentz factor $\gamma$ becomes imaginary or undefined. Reduces to Galilean transformation when $v\ll c$.
\end{itemize}

\item \textbf{Newton's First Law (Inertia):} A body stays at rest or moves in a straight line at constant speed unless a net force acts on it.
\[
\sum\mathbf F=0 \;\Rightarrow\; \mathbf v=\text{const}
\]
\begin{itemize}
\item \textbf{Assumptions:} The reference frame is inertial (not accelerating). The body is isolated from external forces.
\item \textbf{Limitations:} Does not apply in non-inertial (accelerating or rotating) frames without adding fictitious forces. In general relativity, gravity is not a force but a curvature of spacetime, so the concept of ``no net force'' must be reinterpreted.
\item \textbf{Constraints:} Defines what an inertial frame is. Cannot be proven independently---it is essentially the definition of an inertial reference frame.
\end{itemize}

\item \textbf{Newton's Law of Universal Gravitation:} Any two masses attract each other with a force that depends on their masses and the distance between them. This single law explains planetary orbits, tides, and the structure of galaxies.
\[
F=\frac{GMm}{r^2}
\]
\begin{itemize}
\item \textbf{Assumptions:} The masses are point-like or spherically symmetric. The gravitational interaction is instantaneous (no propagation delay). Gravity is always attractive.
\item \textbf{Limitations:} Fails for strong gravitational fields (need general relativity). Does not account for the finite speed of gravity (propagation at $c$). Cannot explain the precession of Mercury's orbit or gravitational lensing.
\item \textbf{Constraints:} Valid in the weak-field, low-velocity regime. For speeds $v\ll c$ and gravitational potentials $\Phi/c^2\ll 1$, it matches general relativity to first order.
\end{itemize}

\item \textbf{Newton's Second Law:} The net force on an object equals its mass times acceleration. This is the central equation of mechanics.
\[
\mathbf F=m\mathbf a
\]
\begin{itemize}
\item \textbf{Assumptions:} The reference frame is inertial. The mass $m$ is constant (no mass change during motion). The force is the net (total) force.
\item \textbf{Limitations:} Does not apply to objects with changing mass (e.g., rockets expelling fuel) without using the more general form $\mathbf F=d\mathbf p/dt$. Fails at relativistic speeds---the correct form is $\mathbf F=d\mathbf p/dt$ where $\mathbf p=\gamma m\mathbf v$.
\item \textbf{Constraints:} Valid for $v\ll c$. In non-inertial frames, fictitious forces must be added. For variable-mass systems, the general form $\mathbf F=d\mathbf p/dt$ must be used.
\end{itemize}

\item \textbf{Newton's Third Law:} Every force has an equal and opposite partner. When you push a wall, the wall pushes back with the same force.
\[
\mathbf F_{12}=-\mathbf F_{21}
\]
\begin{itemize}
\item \textbf{Assumptions:} The forces act along the line joining the two objects (central forces). The system is non-relativistic. No time delay in force propagation.
\item \textbf{Limitations:} Does not hold for magnetic forces between moving charges (the forces are equal and opposite but not necessarily along the line joining the charges). In electrodynamics, the law must be generalized to include field momentum. Fails at relativistic speeds.
\item \textbf{Constraints:} Valid for contact forces and central forces. Conservation of momentum follows from Newton's third law plus Newton's second law.
\end{itemize}

\item \textbf{Schr\"odinger Equation:} The central equation of quantum mechanics. It predicts how the wave function $\psi$ of a particle evolves over time. If you know the wave function, you can calculate the probability of finding the particle anywhere.
\[
i\hbar\frac{\partial\psi}{\partial t}=\hat H\psi
\]
\begin{itemize}
\item \textbf{Assumptions:} The system is non-relativistic (the Hamiltonian contains $p^2/2m$, not the relativistic energy--momentum relation). The wave function $\psi$ provides a complete description of the system. Spin is not included unless Pauli's equation is used.
\item \textbf{Limitations:} Does not incorporate special relativity (the Dirac or Klein--Gordon equations are needed). Cannot describe particle creation or annihilation (need quantum field theory). The many-body problem is intractable analytically for more than a few particles.
\item \textbf{Constraints:} Valid for $v\ll c$. Requires a well-defined potential $V(\mathbf r,t)$. The Hamiltonian must be Hermitian to ensure real energy eigenvalues and unitary time evolution.
\end{itemize}

\end{enumerate}

\section*{What Richard Feynman Would Have Asked}

\subsection*{1. The Plain-English Translation}
\textit{``Don't just repeat the symbols to me. What is this equation actually saying in plain, everyday language?''}

The Core Meaning: These equations are the grammar of the universe. Conservation laws say certain quantities---energy, momentum, charge---cannot be created or destroyed, only moved around. Maxwell's equations say that shaking an electric charge creates a ripple that travels at the speed of light, and that light itself is an electromagnetic wave. Einstein's $E = mc^2$ says that a tiny speck of matter holds a stupendous amount of energy, like a compressed spring waiting to release. The Schr\"odinger equation says that quantum particles behave like waves, with all the interference and fuzziness that implies.

The Metaphor: Think of these equations as the rules of a game. You did not invent the game---nature did. These equations describe what moves are legal. Conservation of energy says you cannot score points out of nothing. Maxwell's equations say that electromagnetic ``plays'' propagate at a fixed speed. The quantum equations say that the game has built-in randomness at its smallest scales.

\subsection*{2. The Localized Mechanics}
\textit{``Look at a tiny, microscopic point in space. What is happening there that these equations describe?''}

He would press you on what each conservation law means at a single point. Conservation of energy at a point means: if energy is increasing somewhere, it must be flowing in from somewhere else---there is no local ``source'' of energy. Conservation of charge at a point means: if charge is accumulating, current must be flowing in. These are continuity equations---the same mathematical structure appears in fluid dynamics, electromagnetism, and quantum mechanics.

The Smoothness Check: He would ask you to explain why Maxwell's equations predict that electromagnetic waves travel at a fixed speed $c = 1/\sqrt{\mu_0 \epsilon_0}$. This is not a coincidence---it is a consequence of the fact that electric and magnetic fields are two aspects of one thing, and their mutual induction creates a self-sustaining wave at a speed determined by the properties of empty space itself.

\subsection*{3. Testing the Extreme Limits}
\textit{``Let's push these equations to their breaking points. Where do they fail?''}

The High-Energy Limit: At very high energies (near the Big Bang), all four forces may have been unified into one. The Standard Model---our best collection of equations---does not include gravity. At the Planck scale ($10^{-35}$ m, $10^{19}$ GeV), spacetime itself may become quantum foam, and all our equations may break down.

The Low-Energy Limit: At everyday energies, quantum effects are negligible and classical equations work beautifully. But at the atomic scale ($10^{-10}$ m), quantum mechanics dominates. The equations do not ``break down'' at these limits---they reveal deeper structure that classical physics misses.

\subsection*{4. Dimensionality and Geometry}
\textit{``Does this equation care about the shape of the space it lives in?''}

He would point out that conservation laws are consequences of symmetries (Noether's theorem). If the laws of physics do not change when you shift time (translation in time symmetry), energy is conserved. If they do not change when you shift space (translation in space symmetry), momentum is conserved. If they do not change when you rotate (rotational symmetry), angular momentum is conserved. The equations we know are not arbitrary---they are dictated by the symmetries of spacetime itself.

\subsection*{5. Cross-Disciplinary Connection}
\textit{``You learned these in physics class, but where else do they appear?''}

Conservation laws appear everywhere: in economics (conservation of money), in ecology (conservation of biomass), in information theory (conservation of information in reversible computing). Maxwell's equations appear in any system where two fields regenerate each other---chemical reaction-diffusion systems, predator-prey populations, and even certain financial models. The deep reason is that the same mathematics describes any system where two quantities mutually sustain each other.


%=============================================================
\part{Individual Equations}
%=============================================================

%=============================================================================
% BATCH 1 — Chapters 1–12
% Physics Equations Reference Book
%=============================================================================

%=============================================================================

